%!TEX root = ../main.tex

\section{Исследование модельных краевых задач}

\subsection{Суть проблемы}

Как говорилось ранее, исследуемая модель \eqref{eq:electrostatic}, \eqref{eq:phi}, в более простой форме предложенная в работе \cite{pitike_dielectric_breakdown}, создана на основе моделей типа диффузной границы из теории трещин. В ряде статей про подобные модели (например, \cite{pitike_dielectric_breakdown}, \cite{sargado_phase_field}) рассматривается краевая задача, приведенная ниже в терминах исследуемой модели канала электрического пробоя.

Рассмотрим задачу с учетом упрощающих допущений в одномерой постановке \eqref{eq:one_dim}. Пусть $\Phi^+ = 0$, $E = 0$ -- электрическое напряжение в системе отсутствует. Будем искать стационарное во времени распределение фазового поля, то есть $\dot{\phi} = 0$. Таким образом, левая часть и первое слагаемое правой части уравнения \eqref{eq:one_dim} равны нулю. Сокращая равенство на ненулевое $\Gamma$, получаем
\begin{equation}
	\cfrac{d^2 \phi}{dx^2} + \cfrac{2}{l^2} f'(\phi) = 0.
	\label{eq:stationary_basic}
\end{equation}
Указанную краевую дифференциальную задачу второго порядка будем рассматривать на луче $[0, +\infty)_x$ с граничными условиями $\phi(0) = 0$, $\phi \to 1$ при $x \to +\infty$.

Домножим обе части равенства на $\phi'_x$. Учитывая, что $f'_\phi \phi'_x = f'_x$ и $2 \phi'_x \phi''_{xx} \hm = [(\phi'_x)^2]'_x$, проинтегрируем уравнение: $(\phi'_x)^2 + (4 / l^2) \cdot f(\phi) = C$. При $x \to +\infty$ согласно граничному условию $\phi \to 1$; естественно также считать, что при этом $\phi'_x \to 0$, следовательно, $C = 4 / l^2$ и
\begin{equation}
	\cfrac{d \phi}{dx} = \cfrac{2}{l} \sqrt{1 - f(\phi)}.
	\label{eq:stationary_basic_transformed}
\end{equation}
Итак, мы перешли к обыкновенной задаче Коши с уравнением \eqref{eq:stationary_basic_transformed} и условием $\phi(0) = 0$, решение которой существует и единственно. Напомним, что $f(\phi) = 4 \phi^3 - 3 \phi^4$, как было введено в начале подраздела \ref{ssec:phase_field_equation}.

Рассмотренная конфигурация системы имеет следующий смысл: найдено распределение фазового поля в полупространстве сбоку от проводящей пластины, состоящей из полностью разрушенного вещества. Связь решения краевой задачи при $E = 0$ с профилем фазового поля в расчете с $E > 0$ и развитием пробоя опишем на уровне интуиции. Автор неявно прибегает к аналогии с задачей теории трещин, на текущем этапе работы пропуская подробности и доказательства ряда утверждений, которые могут в дальнейшем стать предметом отдельного исследования. Итак, рассмотрим одномерную постановку задачи с условием \eqref{cond:boundary_charge} постоянства заряда $q^+ > 0$ на обкладке; пусть в образце развился канал пробоя с центром в координате $x_0$, $\phi(x_0) = \phi_0 \ll 1$. Пусть параметр $\delta$ в формуле \eqref{eq:epsilon} диэлектрической проницаемости $\epsilon(\phi)$ стремится к $0$ (или вовсе равен $0$). При рассмотрении все большего $q^+ \to +\infty$ выполняется $\phi_0 \to +0$, $\epsilon(\phi_0) \to +\infty$,
\[
	\int\limits_0^L \epsilon dx \to +\infty,
\]
но $E \to +0$, то есть за счет сильного увеличения диэлектрической проницаемости образца зависимость $E$ от $q^+$ не прямая, а обратная. Таким образом $E$, постоянное во всем образце в упрощенной постановке задачи, стремится к $0$, и профиль $\phi$ с каждой стороны от точки $x_0$ стремится к решению описанной выше краевой задачи.

Краевая задача \eqref{eq:stationary_basic_transformed} к тому же показывает влияние на систему параметра~$l$: видно, что $l$ в уравнении~\eqref{eq:stationary_basic_transformed} есть коэффициент масштаба решения вдоль оси $x$. Можно показать \cite{zipunova_higher_codimension}, что при $x > l$ решение рассматриваемой задачи Коши $\phi \approx 1$. Другими словами, ее решение есть распределение фазового поля, локализованное на отрезке $[0, l]$.

Подобный анализ вполне подходит для задачи из теории трещин (вместо проводящей пластины была бы плоская трещина). Однако характерный канал пробоя~-- объект не двумерный, а одномерный. Проверим, можно ли провести аналогичное рассуждение не для пластины, а для тонкого прямого проводника.

Как и ранее, электрическое напряжение нулевое: $E = 0$. Рассмотрим задачу в постановке, подобной \eqref{eq:one_dim}, но в цилиндричесской области $\clOmega$, и перейдем в цилиндрические координаты $r$, $\theta$, $y$. Теперь фазовое поле $\phi$ и диэлектрическая проницаемость $\epsilon$ однородны по $\theta$ и $y$ и могут рассматриваться как функции радиуса $r = \sqrt{x^2 + z^2}$. Будем искать решение аналогичной краевой задачи: $\phi|_{r = 0} = 0$, $\phi \to 1$ при $r \to +\infty$. Так как $\phi = \phi(r)$, выражение для лапласиана $\phi$ в цилиндрических координатах принимает вид
\[
	\Delta \phi = \cfrac{1}{r} \partr{} \left( r \partr{\phi} \right) = \cfrac{1}{r} \partr{\phi} + \partrr{\phi}.
\]
С учетом этого цилиндрическая задача описывается уравнением
\begin{equation}
	\cfrac{1}{r} \partr{\phi} + \partrr{\phi} + \cfrac{2}{l^2} f'(\phi) = 0.
	\label{eq:stationary_basic_cylindrical}
\end{equation}

Подобное рассуждение проделано в работе \cite{zipunova_higher_codimension}; за ним следует анализ уравнения~\eqref{eq:stationary_basic_cylindrical}. На основании теоретических результатов из работы \cite{cirstea_elliptic_equations} заключается, что поставленная краевая задача некорректна и решения не имеет. Даже на уровне интуиции постановка задачи выглядит необычно: условие $\phi|_{x,z = 0} = 0$ задано не на двумерной, а на одномерной <<внутренней границе>> области $\Omega$.

Возникает желание формально изменить модель так, чтобы описанная краевая задача имела решение. При моделировании канал пробоя невозможно явно представить сегментом линии, за исключением тривиальных случаев,~-- однако естественно считать его <<нитевидной>> областью соответствующей формы. При $q^+ \to +\infty$, $E \to +0$, как обсуждалось выше, распределение фазового поля вблизи канала пробоя должно приближаться к решению краевой задачи.


\subsection{Предложенное обобщение}

В ответ на описанную в предыдущем подразделе проблему в работе \cite{zipunova_higher_codimension} на основании теоретических результатов работ \cite{sobolev_functional_analysis, oleynik_biharmonic_equations, sternin_elliptic_equations, lewis_quasi_linear} предлагается следующее обобщение уравнения динамики фазового поля, при котором постановка условий на границах размерности~1 (соответственно, коразмерности~2) является математически корректной:
\begin{equation}
	\cfrac{1}{m} \partt{\phi} = \half \epsilon'[\phi] (\vE, \vE) + \cfrac{\Gamma}{l^2} f'(\phi) + \cfrac{\Gamma}{2} \Delta \phi - \alpha \cfrac{\Gamma l^2}{4} \Delta^2 \phi + \beta \Gamma l^{p - 2} \Div \left( \| \nabla \phi \|_2^{p - 2} \nabla \phi \right).
	\label{eq:phi_corrected}
\end{equation}
Здесь $\alpha, \beta \geqslant 0$~-- некоторые константы, $p$~-- четное натуральное число, не меньшее~4. Дифференциальный оператор $\Div (\| \nabla \phi \|_2^{p - 2} \nabla \phi)$ принято называть \emph{$p$-лапласианом}, $\Delta^2 \phi \hm = \Delta(\Delta \phi)$~-- \emph{билапласианом}. В дальнейшем будем считать $p = 4$. Отметим, что слагаемое с $p$-лапласианом, хотя отсутствовало в оригинальной статье \cite{pitike_dielectric_breakdown}, но входило в исходное уравнение \eqref{eq:phi}, однако в упрощающах допущениях постановки \eqref{eq:one_dim} было на время исключено из рассмотрения.